\documentclass[11pt,a4paper]{article}

\usepackage{fontspec}
\setmainfont{Caladea}
\setsansfont{Carlito}
\setmonofont{DejaVu Sans Mono}[Scale=0.82]
\usepackage{unicode-math}
\setmathfont{latinmodern-math.otf}
\usepackage[french]{babel}
\usepackage[a4paper,margin=2.4cm,headheight=14pt]{geometry}
\usepackage{microtype}
\usepackage[locale=FR,group-minimum-digits=4,mode=text]{siunitx}
\usepackage{booktabs,array,tabularx,longtable}
\usepackage{enumitem}
\usepackage{xcolor}
\usepackage{tikz}
\usetikzlibrary{babel,arrows.meta,positioning,calc,fit,backgrounds,decorations.pathreplacing}
\usepackage{pgfplots}
\pgfplotsset{compat=1.18}
\usepgfplotslibrary{fillbetween}
\usepackage[most]{tcolorbox}
\usepackage{fancyhdr}
\usepackage[hidelinks]{hyperref}

\definecolor{encre}{HTML}{1F3A5F}
\definecolor{bleupale}{HTML}{E8EEF6}
\definecolor{ambre}{HTML}{8A5A00}
\definecolor{ambrepale}{HTML}{FBF1DD}
\definecolor{rouge}{HTML}{8E2A2A}
\definecolor{rougepale}{HTML}{F8E6E6}
\definecolor{vert}{HTML}{2E5E2E}
\definecolor{vertpale}{HTML}{E7F1E4}
\definecolor{gris}{HTML}{5A5A5A}

\hypersetup{colorlinks=true,linkcolor=encre,citecolor=encre,urlcolor=encre}

\newtcolorbox{encadre}[1]{enhanced,breakable,colback=bleupale,colframe=encre,
  boxrule=0pt,leftrule=2.5pt,arc=0pt,left=8pt,right=8pt,top=6pt,bottom=6pt,
  fonttitle=\sffamily\bfseries\small,coltitle=encre,colbacktitle=bleupale,
  title={#1},attach title to upper,after title={\par\smallskip}}
\newtcolorbox{alerte}[1]{enhanced,breakable,colback=ambrepale,colframe=ambre,
  boxrule=0pt,leftrule=2.5pt,arc=0pt,left=8pt,right=8pt,top=6pt,bottom=6pt,
  fonttitle=\sffamily\bfseries\small,coltitle=ambre,colbacktitle=ambrepale,
  title={#1},attach title to upper,after title={\par\smallskip}}
\newtcolorbox{correction}[1]{enhanced,breakable,colback=rougepale,colframe=rouge,
  boxrule=0pt,leftrule=2.5pt,arc=0pt,left=8pt,right=8pt,top=6pt,bottom=6pt,
  fonttitle=\sffamily\bfseries\small,coltitle=rouge,colbacktitle=rougepale,
  title={#1},attach title to upper,after title={\par\smallskip}}

\newcommand{\terme}[1]{\emph{#1}}
\newcommand{\unun}{{\NoAutoSpacing 1:1}}
\newcommand{\unN}{{\NoAutoSpacing 1:N}}
\newcommand{\brkunder}{\textunderscore\allowbreak}
\newcommand{\cle}[1]{{\ttfamily\let\_\brkunder #1}}
\newcommand{\nat}[1]{\textsf{\small #1}}
\newcommand{\fort}{$\bullet\bullet$}
\newcommand{\moy}{$\bullet$}
\newcommand{\faible}{$\circ$}

\pagestyle{fancy}
\fancyhf{}
\fancyhead[L]{\sffamily\footnotesize\color{gris}Study 1B · Phase A · Hypothèses d'expérimentation}
\fancyhead[R]{\sffamily\footnotesize\color{gris}v0.4}
\fancyfoot[C]{\sffamily\footnotesize\thepage}
\renewcommand{\headrulewidth}{0.3pt}

\setlist{itemsep=2pt,topsep=4pt}
\setlength{\parskip}{0.45em}
\setlength{\parindent}{0pt}
\setlength{\emergencystretch}{2.5em}

\begin{document}

\begin{titlepage}
\vspace*{3cm}
{\sffamily\color{gris}\small Programme de compression d'embeddings biométriques\par}
\vspace{0.4cm}
{\Huge\bfseries\color{encre}Study 1B — Phase A\par}
\vspace{0.3cm}
{\LARGE Hypothèses d'expérimentation\par}
\vspace{0.3cm}
{\large\itshape Du contrat scientifique à l'expérience exécutable\par}
\vspace{1.6cm}
\begin{tabular}{@{}ll@{}}
\textsf{\small Version} & 0.4 \\
\textsf{\small Rattachement} & head \cle{cb43fdb1543700ae1429a0ed93ea99eed9065899} \\
\textsf{\small Nature} & spécification d'assurance d'ingénierie \\
\textsf{\small Compagnon de} & \cle{STUDY1B\_PHASE\_A\_CONTRACT\_TO\_EXECUTION\_TRACEABILITY\_V0\_4.md} \\
\end{tabular}
\vfill
\begin{encadre}{Statut}
Ce chapitre n'est ni une revue, ni un verdict, ni une autorisation d'exécuter. Il
consigne les hypothèses sur lesquelles repose la Phase A, la manière dont chacune
doit être gardée, et les décisions qui restent à prendre. Aucun résultat biométrique
protégé n'a été ouvert pour l'écrire.

Il a été rédigé avec l'assistance d'un modèle de langage. Les affirmations de ce
modèle ne sont pas des sources : toute valeur factuelle porte sa provenance
(section~\ref{sec:provenance}). Le chapitre documente lui-même une erreur factuelle
commise puis corrigée pendant sa rédaction (section~\ref{sec:correction}).
\end{encadre}
\end{titlepage}

\tableofcontents
\clearpage

% =====================================================================
\section{Objet du chapitre et lecteurs visés}
% =====================================================================

Ce chapitre s'adresse à deux lecteurs qui n'ont pas le même vocabulaire. L'ingénieur
système connaît la mémoire, les caches, la bande passante et les budgets de
puissance, mais pas nécessairement la manière dont on qualifie une dégradation
biométrique. Le scientifique biométricien connaît les taux d'erreur, les points de
fonctionnement et la non-infériorité, mais pas nécessairement pourquoi une recherche
dans une galerie est limitée par la mémoire plutôt que par le calcul. Le
chapitre suppose l'un ou l'autre, jamais les deux. La section~\ref{sec:vocabulaire}
établit un vocabulaire commun, et le glossaire de l'annexe~\ref{sec:glossaire}
définit chaque sigle.

Il répond à une question précise : \emph{sur quelles hypothèses repose la Phase A,
et comment s'assure-t-on que l'expérience exécutée sera bien celle qui a été
conçue ?} Cette seconde moitié est la raison d'être du document. Au terme de sept
cycles de revue, le programme disposait d'un contrat scientifique soigneusement
revu, mais d'aucun mécanisme vérifiant que, le jour de l'exécution, la machine
ferait ce que le contrat demande. Le contrat de mesure, qui porte toute la
réponse scientifique, n'était ouvert par aucun test.

Le chapitre avance en trois temps. Il pose d'abord le modèle physique qui gouverne
la mesure, et l'hypothèse nulle qui en découle. Il reprend ensuite chaque
hypothèse — de mesure, d'appareil, opérationnelle, de décision — en disant ce
qu'elle garantit et ce qu'on conclurait à tort si elle était violée sans que
personne le voie. Il termine par les décisions qui doivent être gelées avant de
choisir la moindre machine, et par les questions qui restent ouvertes.

% =====================================================================
\section{La question posée, et celle qui ne l'est pas}
% =====================================================================

\subsection{D'où vient Study 1B}

Le programme cherche à savoir si l'on peut comprimer une représentation faciale de
512 dimensions en 128 sans la rendre biométriquement inacceptable. Trois études se
succèdent. \textbf{Study 0} comparait quatre routes — la représentation brute en
512 dimensions et trois compressions en 128 : projection aléatoire, analyse en
composantes principales, et projection apprise de type siamois — sur un extracteur
ResNet-18 entraîné sur ImageNet. Après correction de l'estimation d'incertitude, son
résultat était négatif ; mais son substrat n'était pas un vrai modèle de
reconnaissance faciale. \textbf{Study 1A} a qualifié ce substrat manquant : un
extracteur AdaFace IR101, entraîné sur WebFace12M, retenu comme référence 512D.
\textbf{Study 1B} reprend la comparaison sur ce substrat qualifié.

Le protocole confirmatoire de Study 1B fixe une marge de non-infériorité
$\delta_{\mathrm{NI}} = 0{,}03$ sur l'écart de taux de faux rejets, mesuré à un taux
de fausses acceptations de \num{0,01}. Deux campagnes de calibration de puissance,
S4N1 et S4N2, ont établi que ce protocole n'atteint pas la puissance requise de
\num{0,90} : \num{0,869} et \num{0,849} pour le candidat préféré. Toutes deux sont
closes négatives. Elles reposent sur la même population simulée et ne constituent
donc qu'une seule ligne de preuve, examinée par deux procédures d'incertitude
différentes. Les résultats biométriques réels de Study 1B — les étapes dites SCREEN
et TEST — restent scellés.

\subsection{Ce que mesure la Phase A}

La Phase A est un banc d'ingénierie. Elle ne répond à aucune question biométrique.
Elle mesure, sur des vecteurs synthétiques, ce que coûte chaque représentation sur
un appareil en périphérie de réseau — un \terme{edge device}, c'est-à-dire un
appareil qui calcule sur place plutôt que d'envoyer ses données à un serveur
central. Trois questions :

\begin{enumerate}
\item \textbf{Faisabilité mémoire} — la galerie tient-elle en RAM sur un palier
  matériel donné, y compris le cas où le bras 512D ne tient pas ?
\item \textbf{Latence et débit} — pour une recherche exhaustive, quel est l'effet de
  la dimension sur le temps de réponse, le débit, la file d'attente et la saturation ?
\item \textbf{Énergie et autonomie} — mesurée à l'entrée de l'appareil, quelle est
  la consommation, et quelles affirmations d'autonomie en découlent légitimement ?
\end{enumerate}

\subsection{Le pare-feu entre les deux pistes}

La règle d'interprétation n°~9 du contrat de benchmark est courte et décisive :
\emph{aucun résultat protégé de Study 1B ne peut être ouvert pour choisir ou régler
une configuration d'ingénierie}. Elle protège contre une contamination précise. Si
l'on connaissait les écarts biométriques réels avant de choisir le matériel, on
serait tenté de choisir le palier où la compression paraît la plus avantageuse. La
section~\ref{sec:decision} montre que cette règle impose un ordre de gel que la
discussion produit avait failli inverser.

% =====================================================================
\section{Vocabulaire commun}\label{sec:vocabulaire}
% =====================================================================

\subsection{Pour le lecteur biométricien : ce qui gouverne la machine}

\paragraph{Bras.} Terme emprunté aux essais comparatifs, comme le bras de traitement
d'un essai clinique. La Phase A compte exactement deux bras : le bras 512D et le bras
128D. Un bras est la configuration de mesure entière sous une seule condition : même
machine, même logiciel, même nombre de fils d'exécution, même galerie, même charge.
Seule la dimension change. On parle de \terme{bras appariés} parce qu'on les fait
tourner dans des conditions identiques pour attribuer l'écart à la seule variable
qui diffère.

\paragraph{Galerie et gabarit.} La galerie est l'ensemble des représentations
stockées contre lesquelles on compare un visage entrant. Chaque personne y figure
par plusieurs gabarits — trois dans le scénario retenu — c'est-à-dire plusieurs
vecteurs issus de captures différentes.

\paragraph{Working set.} L'ensemble des données qu'un programme touche réellement
pendant son travail. Pour une recherche exhaustive, c'est la galerie entière :
chaque requête compare le vecteur entrant à tous les vecteurs stockés.

\paragraph{Hiérarchie mémoire et cache.} Un processeur ne lit pas directement la
mémoire principale ; il passe par des niveaux de cache de plus en plus grands et de
plus en plus lents — L1, L2, puis L3, appelé aussi \terme{LLC} pour \emph{last level
cache}. Un accès en cache coûte environ un ordre de grandeur de moins qu'un accès en
RAM. Le processeur décide seul de ce qui reste en cache : on ne \og charge \fg{} pas
une galerie en cache, on constate seulement si elle y tient.

\paragraph{Bande passante mémoire.} Le nombre d'octets par seconde que la mémoire
peut livrer au processeur. Pour une recherche exhaustive sur une grande galerie,
c'est elle — et non la vitesse de calcul — qui limite le débit : le processeur passe
l'essentiel de son temps à attendre les données.

\paragraph{Saturation et régime.} Un appareil sature quand la charge demandée
dépasse ce qu'il peut traiter : les requêtes s'accumulent dans une file et le temps
de réponse n'est plus du temps de calcul mais du temps d'attente. On parle de
\terme{régime} pour désigner l'état d'un bras à une charge donnée : non saturé, en
transition, saturé. Comparer un bras saturé à un bras qui ne l'est pas revient à
comparer de l'attente à du service.

\paragraph{Throttling.} Réduction automatique de la fréquence du processeur quand sa
température dépasse un seuil. C'est une marche d'escalier, pas une pente douce :
sous le seuil, pleine vitesse ; au-dessus, vitesse réduite.

\paragraph{Clause, garde, et les trois familles de garde.} Une \terme{clause} dit ce
qui doit être vrai — par exemple, \og le swap doit être nul pendant la mesure \fg.
C'est du texte dans un contrat. Une \terme{garde} est le mécanisme qui vérifie et
qui bloque. Il en existe trois familles, qui n'attrapent pas les mêmes erreurs. Le
\terme{test de contrat} vérifie que le document dit encore ce qu'il disait, et
attrape une modification silencieuse du contrat. L'\terme{assertion d'exécution}
vérifie que la machine a réellement fait ce que le document demande, et attrape un
swap actif ou un wattmètre trop lent. Le \terme{contrôle humain} couvre ce qui ne
s'automatise pas.

\subsection{Pour le lecteur ingénieur : ce qui gouverne la décision biométrique}

\paragraph{Vérification et identification.} En vérification \og un contre un \fg{}
(\unun), on compare une personne à une identité revendiquée, comme au contrôle d'un
titre d'embarquement. En identification \og un contre plusieurs \fg{} (\unN), on
cherche une personne dans toute une galerie, comme dans une liste de surveillance.
Ce ne sont pas les mêmes métriques.

\paragraph{FMR et FNMR.} En \unun, le \terme{FMR} (\emph{false match rate}) est la
proportion d'imposteurs acceptés à tort ; le \terme{FNMR} (\emph{false non-match
rate}) est la proportion de personnes légitimes rejetées à tort. On règle le seuil
de décision pour fixer le FMR, puis on mesure le FNMR qui en résulte : c'est le
\terme{point de fonctionnement}.

\paragraph{FPIR et FNIR.} En \unN, les grandeurs correspondantes sont le \terme{FPIR}
(\emph{false positive identification rate}), proportion de personnes absentes de la
liste qui déclenchent une alerte, et le \terme{FNIR} (\emph{false negative
identification rate}), proportion de personnes de la liste qui passent sans
alerte. Un FMR \unun{} multiplié par la taille de la galerie n'est pas un FPIR.

\paragraph{Non-infériorité.} Démontrer qu'une méthode n'est pas \emph{pire} qu'une
référence au-delà d'une marge fixée d'avance — ici, que le FNMR de la route 128D ne
dépasse pas celui de la référence 512D de plus de trois points. C'est une question
différente de \og est-elle meilleure ? \fg, et elle exige que la marge soit fixée
avant de voir les données.

% =====================================================================
\section{Le modèle de premier ordre et son hypothèse nulle}\label{sec:modele}
% =====================================================================

\subsection{Deux galeries, deux bras, quatre charges utiles}

Le contrat retient deux galeries, marquées toutes deux comme des scénarios de
dimensionnement et non comme des effectifs observés. \textbf{G0} est une liste de
\num{10000} identités, soit \num{30000} vecteurs à trois gabarits par personne.
\textbf{G2} est une galerie centrale de \num{510000} identités, soit
\num{1530000} vecteurs. En virgule flottante sur 32 bits (FP32), chaque composante
occupe 4 octets, et chaque vecteur $4d$ octets.

\begin{table}[h]
\centering\small
\begin{tabular}{@{}lrrr@{}}
\toprule
& vecteurs & bras 512D & bras 128D \\
\midrule
G0 — liste de surveillance & \num{30000} & \num{61,44}~Mo & \num{15,36}~Mo \\
G2 — galerie centrale & \num{1530000} & \num{3133440000} octets & \num{783360000} octets \\
\bottomrule
\end{tabular}
\caption{Charge utile de galerie par bras. Le rapport est exactement 4 par
construction ; c'est une valeur \emph{dérivée}, pas un résultat de mesure.}
\label{tab:charges}
\end{table}

\subsection{Pourquoi la recherche est limitée par la mémoire}

Le cadre qui gouverne tout le raisonnement de la Phase A est le \terme{modèle
Roofline} de Williams, Waterman et Patterson~\cite{roofline}. Il dit qu'un calcul
atteint au mieux le minimum entre deux plafonds : la crête de calcul du processeur,
et le produit de la bande passante mémoire par l'\terme{intensité opérationnelle} —
le nombre d'opérations effectuées par octet lu en mémoire.

\[
  \text{performance atteignable} = \min\bigl(\text{crête de calcul},\;
  \text{bande passante} \times \text{intensité opérationnelle}\bigr)
\]

Une recherche exacte non groupée compare la requête à chaque vecteur de la galerie :
$d$ produits et $d$ sommes, soit $2d$ opérations, pour $4d$ octets lus.
L'intensité opérationnelle vaut donc
\[
  \frac{2d}{4d} = 0{,}5 \ \text{opération par octet},
\]
\textbf{quelle que soit la dimension}. C'est très en deçà du point de crête de
tout processeur usuel : les deux bras sont profondément limités par la mémoire.

\begin{figure}[h]
\centering
\begin{tikzpicture}
\begin{loglogaxis}[
  width=12.5cm,height=7cm,
  xmin=0.1,xmax=100,ymin=0.05,ymax=20,
  xlabel={intensité opérationnelle (opérations par octet)},
  ylabel={performance atteignable},
  xtick={0.1,0.5,1,8,10,100},
  xticklabels={{0,1},{0,5},1,8,10,100},
  ytick=\empty,
  axis lines=left,
  every axis label/.style={font=\small},
  tick label style={font=\small},
]
\addplot[very thick,encre,domain=0.1:12,samples=2] {x};
\addplot[very thick,encre,domain=12:100,samples=2] {12};
\node[font=\footnotesize,encre,rotate=33] at (axis cs:1.2,1.9) {plafond de bande passante};
\node[font=\footnotesize,encre] at (axis cs:40,15.5) {plafond de calcul};
\draw[dashed,gris] (axis cs:12,0.05) -- (axis cs:12,12);
\node[font=\footnotesize,gris,anchor=west] at (axis cs:12.5,0.09) {point de crête};
\addplot[only marks,mark=*,mark size=3pt,rouge] coordinates {(0.5,0.5)};
\node[font=\footnotesize,rouge,anchor=north west,align=left] at (axis cs:0.52,0.42)
  {scan exact non groupé\\les deux bras, même point};
\draw[-{Stealth[length=3mm]},thick,ambre] (axis cs:0.6,0.6) -- (axis cs:7,7);
\addplot[only marks,mark=o,mark size=3pt,ambre,thick] coordinates {(8,8)};
\node[font=\footnotesize,ambre,anchor=south east,align=right] at (axis cs:7,8.5)
  {groupement de $B$ requêtes\\intensité $0{,}5\times B$};
\end{loglogaxis}
\end{tikzpicture}
\caption{Lecture Roofline de la Phase A. Figure schématique : la position du point de
crête dépend de la plateforme et n'est pas une valeur mesurée.}
\label{fig:roofline}
\end{figure}

La figure~\ref{fig:roofline} porte la conséquence la plus importante du chapitre. Les
deux bras occupent \emph{le même point} du diagramme : même intensité, donc même
débit d'opérations. Mais le bras 512D a quatre fois plus d'opérations à faire, et
quatre fois plus d'octets à lire. Il met donc quatre fois plus de temps.

\begin{encadre}{Hypothèse nulle de la Phase A}
Tant que les deux bras restent limités par la mémoire, le modèle prédit un rapport
de latence d'\textbf{exactement 4×}, sur toutes les plateformes et à toutes les
charges non saturées. Tout écart mesuré à 4× est donc de l'\emph{information}, pas
du bruit : il signale un mécanisme que le modèle ignore — résidence en cache, choix
différent de micro-noyau par la bibliothèque de calcul, ou throttling thermique. Les
gardes B1, B3 et D9 existent pour rendre ces écarts attribuables.
\end{encadre}

\subsection{Ordres de grandeur par plateforme}

Au premier ordre, le temps de service d'une requête est la charge utile divisée par
la bande passante effective, et le débit de saturation en est l'inverse. Le
tableau~\ref{tab:saturation} applique ce calcul à trois plateformes edge
documentées. L'efficacité de lecture continue — la fraction de la bande passante
théorique réellement obtenue — est \textbf{supposée} à \SI{70}{\percent} ; elle
n'est pas mesurée, et c'est précisément ce que la caractérisation D4 devra
établir.

\begin{table}[h]
\centering\small
\begin{tabular}{@{}lrrrr@{}}
\toprule
& \multicolumn{2}{c}{bande passante (Go/s)} & \multicolumn{2}{c}{saturation G2 (requêtes/s)} \\
\cmidrule(lr){2-3}\cmidrule(l){4-5}
plateforme & crête & effective & 512D & 128D \\
\midrule
Raspberry Pi 5 & 17 & 11,9 & 3,8 & 15,2 \\
Jetson Orin Nano 4 Go & 34 & 23,8 & \emph{infaisable} & 30,4 \\
Jetson Orin Nano 8 Go & 68 & 47,6 & 15,2 & 60,8 \\
Intel N100 & 38,4 & 26,9 & 8,6 & 34,3 \\
\bottomrule
\end{tabular}
\caption{Modèle de premier ordre, efficacité de \SI{70}{\percent} supposée. La
crête du N100 est dérivée (\num{4800} MT/s × 8 octets, canal unique) ; les autres
sont de fiche. G0 sature entre environ 190 et \num{3100} requêtes par seconde selon
le cas, très au-delà de la grille de charge.}
\label{tab:saturation}
\end{table}

Le tableau fait déjà apparaître deux constats que la suite exploite. Sur le
Raspberry Pi 5, le bras 512D sature \emph{à l'intérieur} de la grille de charge
prévue, qui va de \num{0,25} à 8 requêtes par seconde, alors que le bras 128D non :
c'est exactement la situation que la garde B2 doit encadrer
(section~\ref{sec:saturation}). Sur le Jetson Orin Nano 8 Go, en revanche, aucun des
deux bras ne sature avant 8 requêtes par seconde : la grille est trop basse pour ce
matériel. Le même protocole produit donc des cellules interprétables différentes
selon la machine — raison pour laquelle la règle d'affectation de régime doit être
gelée \emph{avant} de choisir la machine.

% =====================================================================
\section{Ce que la vérification a corrigé}\label{sec:correction}
% =====================================================================

\begin{correction}{Correction d'une affirmation fausse}
Une version antérieure de ce travail affirmait que les caches de dernier niveau des
processeurs edge vont \og de 4 à \num{32}~Mo \fg, et en déduisait que la
galerie G0 — \num{61,44}~Mo contre \num{15,36}~Mo — encadrait cette
gamme : le bras 128D y aurait tenu en cache, le 512D non. \textbf{C'était faux.}
L'affirmation venait de l'apprentissage du modèle de langage, non d'une source.
\end{correction}

Les valeurs documentées sont nettement plus basses : \num{2}~Mo de L3 partagé
sur le Raspberry Pi 5~\cite{rpidoc}, \num{2}~Mo par grappe de cœurs sur le
Jetson Orin Nano~\cite{ridgerun}, \num{6}~Mo sur l'Intel N100~\cite{nbcheck}.
Sur aucune de ces trois plateformes, aucun bras de G0 ne tient en cache. Les quatre
combinaisons bras × galerie sont en mémoire principale.

La conséquence n'est pas cosmétique. G0 était décrite comme un
\emph{contrôle à petit working set}, censé isoler un régime où l'arithmétique et le
cache dominent. Sur du matériel edge réaliste, G0 diffère de G2 par la grandeur, pas
par la nature. Un véritable contrôle résident en cache demanderait environ
\num{4000} vecteurs à 128D, ou \num{1000} à 512D, pour \num{2}~Mo de cache.

Le Jetson a par ailleurs validé dans les faits une règle posée auparavant de façon
abstraite : son cache L3 est partagé par grappe, non par puce. Un banc épinglé sur
la grappe de quatre cœurs dispose de \num{2}~Mo, pas de la somme. C'est ce
que la décision D3 appelle le cache \emph{effectif}.

Une seconde leçon est venue en cherchant la fiche Intel du N100. Le premier résultat
était un résumé généré par intelligence artificielle d'une page Intel archivée ; il
attribuait au N100 le double et le triple canal mémoire, l'Hyper-Threading et la
mémoire ECC — trois affirmations que toutes les autres sources contredisent. Un
résumé automatique d'une source primaire n'hérite pas de son autorité. Il est exclu
des sources de ce chapitre, au même titre que les réponses d'un assistant, y
compris celles qui ont servi à le rédiger.

% =====================================================================
\section{Hypothèses de mesure}\label{sec:mesure}
% =====================================================================

Chaque hypothèse est présentée par ce qu'elle garantit, puis par ce qu'on conclurait
à tort si elle était violée sans que personne le voie. Cette seconde question est
celle qui ordonne les priorités : toutes les clauses ne méritent pas le même effort
d'ingénierie, et c'est la gravité de la conclusion fausse qui décide.

\subsection{Faisabilité mémoire}

\paragraph{Résidence en RAM et absence de swap.} Le contrat exige que la galerie
reste entièrement en mémoire vive et que le système n'écrive rien sur disque pour
compenser. Violée, cette clause ferait conclure que le bras 512D tient en RAM sur un
palier alors que le système d'exploitation l'a sauvé en débordant sur le stockage.
C'est la conclusion la plus coûteuse du programme : elle validerait un déploiement
matériellement infaisable. Garde requise : assertion d'exécution sur les compteurs de
swap, avant et après chaque fenêtre de mesure.

\paragraph{Infaisabilité comme résultat.} Si un bras ne tient pas, le contrat
enregistre \cle{RAM\_INFEASIBLE\_ON\_TIER} comme un résultat valide, et interdit
toutes les échappatoires : réduire la galerie, changer le type numérique, activer le
swap, ou migrer un seul bras vers un palier supérieur. Violée, un contournement
passerait pour une mesure appariée.

\paragraph{RAM requise totale.} La faisabilité ne se prononce pas sur les octets de
galerie mais sur la somme de la galerie, du processus, de l'environnement
d'exécution, du système, des tampons, des autres services et d'une marge de
sécurité. Sans marge gelée à l'avance, un bras déclaré faisable avec \SI{2}{\percent}
de mémoire libre ne l'est pas en production — et une marge fixée après avoir vu les
chiffres sera fixée à ce qui fait passer la conclusion voulue.

\paragraph{Cache et bande passante — garde B1.} L'objectif affiché du contrat
distingue un régime où l'arithmétique domine et un régime où la bande passante
domine. Sans la taille du cache et une bande passante mesurée, cette distinction
n'est pas vérifiable. La garde B1 exige de les enregistrer, avec un rapport
\emph{dérivé} et toujours calculable : charge utile × débit complété ÷ bande
passante mesurée. Ce rapport a une propriété utile : il suppose que tout vient de la
RAM, donc il est une borne haute du trafic réel. Une valeur supérieure à 1 n'est pas
une erreur de mesure, c'est la preuve qu'une partie du working set réside en cache.

\subsection{Latence, débit et saturation}\label{sec:saturation}

\paragraph{Sémantique de recherche.} Recherche exacte, un seul meilleur candidat,
aucun seuil, aucun élagage, aucun index. Cette construction fixe le travail par
requête et neutralise un confondant subtil : entre vecteurs aléatoires, les scores de
similarité se concentrent en $1/\sqrt{d}$, deux fois plus dispersés à 128D qu'à 512D.
Tout mécanisme dépendant d'un seuil produirait plus de candidats pour le bras 128D,
et l'écart de dimension serait un artefact du générateur de données.

\paragraph{Isolement de l'intervention.} Matériel, type numérique, cardinalité,
threads et bibliothèque identiques, et surtout : la matrice 512D ne doit pas rester
résidente pendant la mesure 128D. Sinon le bras 128D serait mesuré sur un cache
chaud et une mémoire déjà peuplée, et paraîtrait meilleur pour une raison sans
rapport avec la dimension.

\paragraph{Génération synthétique.} Graine fixée, loi normale en FP32, normalisation,
rejet des valeurs non finies. La marge est confortable : la composante typique vaut
\num{0,0442} à 512D et \num{0,0884} à 128D, très loin des valeurs dénormalisées qui
ralentiraient le calcul.

\paragraph{Génération de charge.} Générateur hors de l'appareil mesuré, arrivées en
boucle ouverte — c'est-à-dire indépendantes des réponses du système —, grille de
\num{0,25} à 8 requêtes par seconde, rafale à 16. Un générateur sur l'appareil
injecterait son propre calcul et sa propre consommation dans la mesure ; une boucle
fermée masquerait la saturation.

\paragraph{Statistiques de mesure.} Cinq redémarrages indépendants, 60~secondes
d'échauffement exclues, fenêtre de 600~secondes, ordre des bras alterné, conservation
des échantillons bruts, aucune exclusion silencieuse de run raté. Chaque élément
ferme une erreur différente : un échauffement non exclu attribue le coût d'un cache
froid à la dimension, un ordre fixe attribue la dérive thermique au second bras, une
exclusion silencieuse installe un biais de survie.

\paragraph{Support des percentiles.} Un percentile 95 exige au moins 200 requêtes
complétées par redémarrage, un percentile 99 au moins \num{1000}. Sur 600~secondes,
la grille donne 150, 300, 600, \num{1200}, \num{2400} et \num{4800} requêtes attendues :
le percentile 95 est supprimé par conception à \num{0,25} requête par seconde, le
percentile 99 jusqu'à 1. Mais une cellule vide a deux sens opposés — charge trop
faible par conception, ou bras saturé qui n'a pas complété assez de requêtes. Publier
les comptes attendus \emph{et} réalisés lève l'ambiguïté.

\paragraph{Régime de saturation — garde B2.} Un rapport 512D/128D n'est publiable que
si les deux bras sont dans le même régime et qu'aucun n'est censuré. La
figure~\ref{fig:regimes} l'illustre sur le Raspberry Pi 5.

\begin{figure}[h]
\centering
\begin{tikzpicture}
\begin{axis}[
  width=12.5cm,height=6.6cm,
  xmin=0,xmax=18,ymin=0,ymax=18,
  xlabel={charge offerte (requêtes/s)},
  ylabel={débit complété (requêtes/s)},
  axis lines=left,
  xtick={0,2,4,8,16},
  ytick={0,3.8,15.2},
  yticklabels={0,{3,8},{15,2}},
  every axis label/.style={font=\small},
  tick label style={font=\small},
  legend style={font=\footnotesize,at={(0.02,0.98)},anchor=north west,draw=none,fill=white},
]
\fill[ambrepale] (axis cs:3.88,0) rectangle (axis cs:18,18);
\node[font=\footnotesize,ambre,align=center] at (axis cs:12.3,7.6)
  {paires à régimes croisés\\aucun rapport publiable};
\addplot[very thick,encre,domain=0:18,samples=200] {min(x,15.2)};
\addlegendentry{bras 128D}
\addplot[very thick,rouge,domain=0:18,samples=200] {min(x,3.8)};
\addlegendentry{bras 512D}
\addplot[dotted,gris,domain=0:18] {x};
\pgfplotsinvokeforeach{0.25,0.5,1,2,4,8,16}{\draw[gris!60] (axis cs:#1,0) -- (axis cs:#1,0.6);}
\end{axis}
\end{tikzpicture}
\caption{Illustration au premier ordre sur Raspberry Pi 5 (efficacité supposée
\SI{70}{\percent}). Les traits courts sur l'axe marquent la grille de charge. À 4
requêtes par seconde le bras 512D est en transition, à 8 et 16 il est saturé, alors
que le bras 128D ne l'est qu'à peine à 16. Seules les charges de \num{0,25} à 2
donnent des paires comparables.}
\label{fig:regimes}
\end{figure}

La garde B2 est aujourd'hui incomplète. Ses cinq étiquettes de régime sont gelées,
mais \emph{la règle qui affecte une cellule à une étiquette ne l'est pas}. Sans elle,
c'est un analyste regardant les données qui décidera après coup si un rapport peut
être publié — exactement la liberté que le programme a supprimée partout ailleurs.
La décision D1 y pourvoit (section~\ref{sec:decisions}).

\paragraph{Régime thermique.} Le throttling est une marche d'escalier. À G2, le bras
512D fait quatre fois plus de trafic mémoire et chauffera davantage. S'il franchit
le seuil et que le bras 128D ne le franchit pas, le rapport mesuré confond
\og quatre fois plus de travail \fg{} et \og à une fréquence plus basse \fg. C'est
l'analogue exact de la saturation, et aucune clause ne le couvrait : il appelle la
même construction, décision D9.

\paragraph{Choix de micro-noyau — garde B3.} La règle d'appariement fixe la
\emph{configuration} de la bibliothèque de calcul, pas le \emph{chemin de code} : une
bibliothèque BLAS choisit couramment des noyaux différents à 512 et à 128
dimensions. Cette différence fait légitimement partie de l'effet mesuré, et la
garde B3 ne cherche pas à l'égaliser ; elle l'enregistre, pour qu'un écart
surprenant puisse être diagnostiqué.

\subsection{Énergie et autonomie}

\paragraph{Plan de mesure.} La mesure canonique se fait à l'entrée de l'appareil ou
à la prise, au moins dix fois par seconde ; la télémétrie interne ne sert que de
diagnostic. Ce choix n'est pas une précaution de principe. Les compteurs intégrés
excluent souvent la mémoire vive, or c'est exactement là que se voit l'écart entre
\num{2,92}~Gio et \num{0,73}~Gio. Une mesure par télémétrie biaiserait
la comparaison en faveur du bras évalué.

\paragraph{Deux conventions d'énergie par identification.} Le total, ligne de base au
repos incluse, est canonique ; l'incrémental, ligne de base soustraite, n'est que
diagnostique, et ses valeurs négatives ne sont pas ramenées à zéro. Écrêter les
négatifs transformerait un bruit symétrique en biais positif. Le total, lui, dépend
fortement de la charge : un chiffre publié sans son point de charge n'a pas de sens.

\paragraph{Modèle de batterie.} L'autonomie est la capacité utile divisée par la
puissance moyenne \emph{totale}, sous un profil de charge nommé. L'énergie
incrémentale ne doit jamais servir à calculer une autonomie : elle la surestimerait
du rapport entre puissance totale et puissance incrémentale, l'erreur la plus facile
à commettre et la plus visible dans un dossier de décision.

\paragraph{Représentativité thermique du banc.} Le contrat enregistre le mode de
refroidissement et la température ambiante, mais n'exige pas qu'ils correspondent au
déploiement. La Fondation Raspberry Pi documente que ses cartes réduisent leur
fréquence dès \SI{80}{\celsius}, davantage à \SI{85}{\celsius}, et que sans
refroidissement un Pi 5 sous charge soutenue reste en throttling permanent — mesuré
à l'air libre sur un banc de laboratoire~\cite{rpiheat}. Un boîtier étanche au
soleil est pire. Mesurer sur banc ventilé et déployer en boîtier sans ventilateur
donne un résultat qui ne se transfère pas, et qui se trompe dans le sens flatteur.

\begin{alerte}{Un avantage candidat de 128D, invisible sur banc ventilé}
À charge égale, le bras 128D fait quatre fois moins de trafic mémoire, donc dissipe
moins. Il peut rester sous le seuil de throttling là où le bras 512D le franchit. Ce
serait un gain opérationnel indépendant de la mémoire vive. Il n'apparaît qu'en
conditions thermiques de déploiement : ventilés, les deux bras restent froids et
l'avantage disparaît de la mesure.
\end{alerte}

% =====================================================================
\section{L'appareil complet}\label{sec:appareil}
% =====================================================================

\subsection{La recherche n'est qu'un consommateur parmi d'autres}

La Phase A mesure la recherche \emph{seule} : le contrat exige qu'aucun autre
processus significatif ne tourne pendant la mesure. C'est la bonne première mesure,
parce qu'elle isole la demande propre de la recherche. Mais un appareil de
reconnaissance déployé à un portique fait bien d'autres choses, qui puisent dans les
mêmes réserves. Quatre budgets sont partagés : la \textbf{puissance}, la
\textbf{chaleur} que le boîtier doit évacuer, la \textbf{bande passante mémoire} et
la \textbf{RAM}. La faisabilité ne se décide donc pas sur le budget total, mais sur le
budget \emph{résiduel}, une fois servis les autres composants.

\begin{figure}[h]
\centering
\begin{tikzpicture}[x=1cm,y=1cm,font=\sffamily\footnotesize]
\def\H{0.9}
\draw[encre,thick] (0,0) rectangle (14,\H);
\node[anchor=south west,encre] at (0,\H+0.08) {budget de puissance de l'appareil — par exemple le plafond PoE 802.3af de 12,95~W};
\fill[gris!25] (0,0) rectangle (2.3,\H);   \node at (1.15,\H/2) {repos};
\fill[gris!40] (2.3,0) rectangle (4.6,\H); \node at (3.45,\H/2) {caméras};
\fill[ambre!35] (4.6,0) rectangle (7.8,\H); \node at (6.2,\H/2) {illumination IR};
\fill[gris!45] (7.8,0) rectangle (10.3,\H); \node at (9.05,\H/2) {capture};
\fill[gris!30] (10.3,0) rectangle (11.4,\H); \node at (10.85,\H/2) {radio};
\fill[encre!85] (11.4,0) rectangle (14,\H); \node[white] at (12.7,\H/2) {recherche};
\draw[decorate,decoration={brace,amplitude=5pt,mirror},encre] (11.4,-0.1) -- (14,-0.1)
  node[midway,below=7pt,align=center,encre] {ce que mesure la Phase A\\et ce qui décide de la faisabilité};
\draw[decorate,decoration={brace,amplitude=5pt,mirror},gris] (0,-0.1) -- (11.3,-0.1)
  node[midway,below=7pt,gris] {autres composants — aucune valeur mesurée à ce jour};
\end{tikzpicture}
\caption{Principe du budget résiduel. Les proportions sont purement illustratives :
aucune n'est mesurée. Le même raisonnement vaut pour la chaleur, la bande passante
et la RAM.}
\label{fig:budget}
\end{figure}

La recherche n'est peut-être pas le plus gros poste. Une illumination infrarouge peut
consommer davantage que le processeur. Et un exemple documenté de caméras extérieures
équipées de chauffage et d'essuie-glace fait état de 32~W par caméra par temps froid
— davantage, à lui seul, que le plafond PoE 802.3af~\cite{ampcom}. Si caméras,
illumination et détection consomment déjà l'essentiel d'un plafond PoE, les quelques
watts d'écart entre les deux bras de recherche sont peut-être précisément ce qui
décide de la faisabilité ; ou ils n'y changent rien. On ne le saura qu'avec les
autres lignes du budget.

\subsection{Les composants}

\begin{table}[h]
\centering\small
\begin{tabularx}{\textwidth}{@{}>{\raggedright\arraybackslash}Xcccc>{\raggedright\arraybackslash}p{4.2cm}@{}}
\toprule
composant & puissance & chaleur & bande passante & RAM & point de vigilance \\
\midrule
calcul — recherche & \multicolumn{4}{c}{\emph{mesuré par la Phase A}} & objet du banc \\
caméra visible & \moy{} & \moy{} & \moy{} & \moy{} & résolution × cadence \\
caméra proche infrarouge (NIR) & \moy{} & \moy{} & \moy{} & \moy{} & sert aussi à la détection de fraude \\
illumination IR & \fort{} & \fort{} & — & — & peut dépasser le calcul ; sécurité oculaire \\
détection de mouvement & \faible{} & \faible{} & — & — & commande la mise en veille \\
chaîne de capture & \fort{} & \fort{} & \fort{} & \fort{} & tourne en continu, rivalise avec la recherche \\
radio & \moy{} & \moy{} & — & \faible{} & commande la propagation de galerie \\
stockage de la galerie & \faible{} & \faible{} & — & — & chiffrement, temps de chargement \\
élément sécurisé & \faible{} & \faible{} & — & — & vol d'appareil, fuite de liste \\
alimentation & pertes & \moy{} & — & — & plafond dur (PoE ou batterie) \\
chauffage & \fort{} & — & — & — & peut dominer l'hiver \\
boîtier & — & décisif & — & — & reflets IR dans la caméra \\
\bottomrule
\end{tabularx}
\caption{Contribution des composants aux quatre budgets. \fort{} charge potentiellement
dominante, \moy{} charge significative, \faible{} charge faible. La chaîne de capture comprend la
détection, l'alignement et l'encodage du visage.}
\label{tab:composants}
\end{table}

\subsection{Quatre interactions qui font dépendre un coût de la taille de galerie}

Ces quatre points relient la liste de composants à l'étude. Dans chacun, un choix de
composant rend un coût opérationnel proportionnel à la taille de la galerie — donc à
la dimension.

\paragraph{Veille et galerie résidente.} Un appareil réveillé par détection de
mouvement dort entre deux passages. Au réveil, la galerie doit être en mémoire. Deux
voies, qui dépendent toutes deux de la taille : soit la RAM reste alimentée en veille,
et ce coût croît avec la capacité ; soit la galerie est rechargée depuis le stockage
et déchiffrée — \num{3,13}~Go contre 783~Mo. Et le réveil n'est utile que si le délai
entre le réveil et la disponibilité est plus court que le temps qu'une personne met à
entrer dans le champ. C'est une latence qui dépend directement de la dimension, et la
Phase A ne la mesure pas.

\paragraph{Radio en rafale et propagation.} Une radio qui ne s'allume que sur alerte
économise l'énergie, mais la mise à jour de la liste de surveillance doit alors
attendre une fenêtre radio. La durée de propagation devient le produit de la taille
de galerie et du cycle radio.

\paragraph{Galerie chiffrée au repos.} Des données biométriques sur un appareil en
espace public doivent être chiffrées sur le stockage. Le déchiffrement au démarrage
ou au réveil coûte un temps et un calcul proportionnels à la taille, et la galerie en
clair occupe quatre fois plus de RAM à 512D.

\paragraph{Bande passante partagée.} La détection de visage et le traitement de
l'image consomment de la bande passante en continu. Or la recherche en est limitée.
Le modèle de premier ordre doit donc utiliser la bande passante \emph{disponible},
pas la bande passante totale : à 17~Go/s sur un Raspberry Pi 5, quelques Go/s pris
par la détection déplacent sensiblement le point de saturation.

% =====================================================================
\section{Hypothèses opérationnelles}\label{sec:operationnel}
% =====================================================================

Les hypothèses opérationnelles relient la mesure au terrain. Le programme en a
aujourd'hui la \emph{structure} complète ; il n'en a pas encore les \emph{valeurs},
qui relèvent du scénario et doivent être déclarées par le responsable du programme.
La section~\ref{sec:questions} en dresse la liste.

\subsection{Classes de déploiement}

\og Dans et autour des stades \fg{} désigne au moins deux situations, et elles ne sont
pas contraintes par la même chose. Un \textbf{portique fixe} dispose vraisemblablement
d'une alimentation par le câble réseau (PoE) ou du secteur : la batterie n'y sert que
de secours, et c'est la chaleur, puis le plafond PoE, qui mordent. La norme 802.3af
garantit \num{12,95}~W à l'appareil alimenté, la 802.3at en garantit \num{25,5}~W — et
ce plafond couvre l'appareil entier. Un \textbf{point mobile ou temporaire} n'a pas
d'infrastructure : sa batterie contraint autant que la chaleur.

La batterie partage d'ailleurs le boîtier avec le processeur. Une cellule lithium-ion
se charge typiquement entre 0 et \SI{45}{\celsius} et se décharge entre $-20$ et
\SI{60}{\celsius} ; la charge sous le point de congélation n'est pas permise~\cite{bu410}.
Le plafond haut compte autant que le bas : dans un boîtier fermé au soleil, chauffé
par le processeur, la cellule peut sortir de sa fenêtre de charge \emph{par le haut}.
Un point mobile en été peut alors fonctionner sans pouvoir se recharger. Un piège de
lecture accompagne ce point : la plage \og operating temperature \fg{} d'une fiche de
cellule reflète souvent la décharge ; la plage de charge est plus étroite~\cite{sunlith}.

\subsection{Du passage à la requête}

La grille de charge est exprimée en requêtes par seconde ; le scénario, en personnes
qui passent. La conversion entre les deux n'est écrite nulle part. Or une personne
qui franchit un portique produit plusieurs images. Chercher chaque visage détecté ou
seulement la meilleure image change la charge d'un facteur 5 à 10 pour le même flux
humain. Et un processeur qui sert plusieurs caméras la multiplie encore : la fiche du
Jetson Orin Nano indique jusqu'à 16 canaux vidéo virtuels~\cite{nvds}.

\subsection{Profils de service}

L'autonomie à \num{0,25} requête par seconde et à 8 diffère bien davantage que tout
écart entre 512D et 128D. Le contrat exige que toute autonomie nomme son profil, mais
aucun profil n'est déclaré. Il en faut deux ou trois, décrits heure par heure : pointe
de flux avant le match, régime de croisière, veille. C'est une décision de scénario,
pas d'analyse : on sait déjà à quoi ressemble une journée de stade. Choisi après coup,
le profil serait choisi pour que la conclusion tombe juste.

\subsection{Conséquence opérationnelle}

\paragraph{Deux familles de métriques, deux branches.} Le contrôle d'accès de type
embarquement se chiffre en FMR et FNMR ; la liste de surveillance en FPIR et FNIR. Les
deux branches de l'analyse de conséquence doivent être construites séparément, et
aucun FMR \unun{} ne doit être multiplié par une taille de galerie pour produire un
risque \unN. Cette erreur classique produit un nombre d'alertes qui paraît rigoureux
et qui est faux d'un facteur inconnu.

\paragraph{Corrélation entre observations.} Une personne vue plusieurs fois — deux
portiques, plusieurs caméras — n'offre pas des chances indépendantes. Si elle est
difficile à reconnaître, elle l'est souvent partout : même visage, mêmes lunettes,
même angle, même éclairage.

\begin{figure}[h]
\centering
\begin{tikzpicture}
\begin{semilogyaxis}[
  width=11cm,height=6cm,
  xmin=1,xmax=5,ymin=5e-6,ymax=0.3,
  xtick={1,2,3,4,5},
  xlabel={nombre d'observations $m$ de la même personne},
  ylabel={probabilité de la manquer partout},
  axis lines=left,
  every axis label/.style={font=\small},
  tick label style={font=\small},
  legend style={font=\footnotesize,at={(0.98,0.98)},anchor=north east,draw=none,fill=white},
]
\addplot[name path=hi,very thick,rouge,domain=1:5,samples=2] {0.1};
\addlegendentry{corrélation forte : $\approx p$}
\addplot[name path=lo,very thick,encre,domain=1:5,samples=50] {0.1^x};
\addlegendentry{indépendance : $p^m$}
\addplot[ambrepale] fill between[of=hi and lo];
\node[font=\footnotesize,ambre,align=center] at (axis cs:3.6,0.004) {bande à publier\\au lieu d'un point};
\end{semilogyaxis}
\end{tikzpicture}
\caption{Pour une probabilité de manquer $p = 0{,}1$ par observation : à trois
observations, \num{0,1}~\% sous indépendance contre \SI{10}{\percent} sous corrélation
forte. Deux ordres de grandeur.}
\label{fig:correlation}
\end{figure}

La règle qui en découle est simple : tout chiffre opérationnel multi-passage ou
multi-appareil se publie comme une \emph{bande} entre ces deux extrêmes, jamais comme
une valeur ponctuelle, sauf à disposer d'un modèle de corrélation justifié.

\begin{encadre}{L'erreur de Study 0, remontée d'un niveau}
Study 0 avait rééchantillonné des paires de comparaisons comme si elles étaient
indépendantes, alors que les mêmes identités y revenaient ; ses intervalles étaient
trop étroits. Traiter comme indépendantes les observations successives d'une même
personne difficile, c'est commettre la même faute un niveau plus haut, avec la même
conséquence : une confiance excessive.
\end{encadre}

\paragraph{Ancrage des références.} Toute valeur de référence nomme sa source, son
jeu de données et son point de fonctionnement, et n'est jamais transportée vers un
autre. Le jeu LFW, par exemple, ne permet pas d'estimer un FMR de $10^{-6}$. Le plafond
ne vient pas du nombre de paires — l'appariement exhaustif en produit de l'ordre de
$10^{8}$ — mais des \num{5749} identités : les paires se recouvrent, et l'information
effective est bornée par le nombre de sujets. C'est encore le même fait statistique
que ci-dessus.

\subsection{Paliers matériels}

L'argument économique le plus fort du dossier est que 128D pourrait éviter le palier
matériel supérieur : un appareil à 4~Go au lieu de 8~Go, par exemple. Trois conditions
le rendent défendable. L'échelle des paliers candidats doit être déclarée \emph{avant}
toute mesure, sans quoi le second palier serait choisi après avoir vu le premier
échouer. La RAM requise doit inclure une marge de sécurité gelée d'avance. Et aucune
conséquence en aval — puissance, batterie, chaleur, masse, coût — ne doit être
inférée du seul écart de RAM : chaque maillon est mesuré ou déclaré supposé.

Une quatrième condition est apparue à la vérification des fiches. Sur le Jetson Orin
Nano, la version 4~Go a une bande passante de crête de 34~Go/s, la version 8~Go de
68~Go/s~\cite{nvds}. Descendre au palier inférieur ne retire pas seulement de la
mémoire : \textbf{il divise la bande passante par deux}. Or la recherche en est
limitée. L'avantage de palier doit donc être évalué avec la bande passante du palier
visé, pas avec celle du palier supérieur.

% =====================================================================
\section{Hypothèses de décision}\label{sec:decision}
% =====================================================================

\subsection{La marge ne bouge pas}

La marge $\delta_{\mathrm{NI}} = 0{,}03$ est une tolérance de recherche gelée, reprise
de Study 0 pour préserver la comparabilité ; elle n'a été fixée ni par un organisme de
normalisation ni par un consensus d'experts. Elle ne peut pas être remplacée à
l'intérieur de Study 1B. Sinon un manque de puissance se résoudrait par un
élargissement de marge, et un résultat \og non démontré \fg{} deviendrait \og démontré \fg{}
par simple redéfinition — la forme la plus grave de raisonnement a posteriori, parce
qu'elle ne laisse aucune trace dans les résultats, seulement dans l'historique du
contrat. Toute marge issue d'un travail de justification ouvre une nouvelle identité
de protocole.

\subsection{L'expert ne voit que des conséquences}

La marge peut en revanche être \emph{justifiée} pour une étude future, par l'avis
d'un expert sur la gravité opérationnelle des dégradations. Le risque est alors que
l'expert, connaissant les écarts réels, trouve acceptable ce qui a été observé. Le
protocole d'élicitation repose sur un principe qui rend presque tout le reste
superflu : \textbf{l'expert ne voit jamais un écart de taux ; il voit des
conséquences}.

\og Plus trois points de FNMR \fg{} est un objet statistique sur lequel un expert n'a
pas d'intuition calibrée, et qui invite à s'ancrer sur 0,03. \og \num{22500} passages
supplémentaires en contrôle manuel sur \num{750000} \fg{} est un objet opérationnel sur
lequel il en a une. Traduire l'écart en conséquence est le travail de l'enveloppe ;
juger l'acceptabilité de la conséquence est celui de l'expert. Séparés ainsi, l'aveuglement
devient structurel : l'expert ne peut pas reconnaître où tombent les valeurs réelles,
parce qu'il ne voit pas l'échelle sur laquelle elles tombent.

\begin{table}[h]
\centering\small
\begin{tabular}{@{}ll@{}}
\toprule
rôle & interdit \\
\midrule
détenteur des résultats SCREEN & construire l'enveloppe, assister à l'élicitation \\
constructeur de l'enveloppe & avoir vu un résultat SCREEN \\
expert & voir un écart de taux, voir un nom de méthode \\
\bottomrule
\end{tabular}
\caption{Séparation des rôles. Si une même personne en tient deux, aucune procédure
ne rétablit l'aveuglement.}
\label{tab:roles}
\end{table}

La grille présentée déborde le plausible des deux côtés, l'ordre de présentation est
aléatoire, chaque cellule montre la bande de corrélation plutôt qu'un point, et les
deux branches — contrôle d'accès et liste de surveillance — sont élicitées séparément.
Une seconde séance, avec des scénarios intercalés reprenant des niveaux déjà jugés
sous une autre forme, vérifie que le seuil énoncé est reproductible ; s'il ne l'est
pas, ce n'est pas une mesure. Le produit n'est pas \og la marge vaut X \fg, mais un
tableau conséquence → acceptable, limite ou inacceptable, par branche. Il donnera
vraisemblablement une marge par branche, et non une marge unique.

\subsection{SCREEN ne s'ouvre qu'après la machine}

La discussion produit a proposé d'ouvrir SCREEN — l'étape exploratoire des résultats
biométriques réels — avant la Phase A, pour disposer plus tôt des écarts réels des
méthodes. Cet ordre violerait la règle n°~9 au sens littéral. Le palier matériel, la
bibliothèque de calcul, le nombre de fils d'exécution et le wattmètre restent à choisir ;
choisis après SCREEN, ils le seraient par des personnes connaissant les écarts
biométriques. La correction ne coûte rien : il suffit d'échanger la position de SCREEN
et de la matérialisation de la plateforme (figure~\ref{fig:gel}).

% =====================================================================
\section{Trois hypothèses oubliées et une question produit}\label{sec:oubliees}
% =====================================================================

Ces trois hypothèses dimensionnent le résultat autant que toutes les précédentes, et
aucune n'est fixée par le contrat. Aucune fiche technique ne les aurait révélées :
elles sont apparues en confrontant le contrat à un cadre — Roofline, les règles de
MLPerf, la théorie des files d'attente — et en posant une seule question :
\emph{qu'est-ce qui, dans ce contrat, n'est fixé nulle part ?} La première a échappé à
la revue de validité de construit, y compris à celle de l'assistant qui a contribué à
ce document.

\subsection{Le groupement des requêtes}

Le contrat fixe l'algorithme, le meilleur candidat unique, l'absence de seuil et
d'index. Il ne dit rien du \terme{batching} — le groupement de plusieurs requêtes
traitées en une seule passe sur la galerie. Le mot n'y apparaît pas.

Or c'est une variable de premier ordre. Grouper $B$ requêtes lit chaque vecteur de la
galerie une fois pour $B$ usages : l'intensité opérationnelle passe de \num{0,5} à
$0{,}5 \times B$ (figure~\ref{fig:roofline}). Tant que l'on reste sous le point de
crête, le débit croît d'environ $B$ fois, pour les deux bras. Une variable non gelée
peut donc déplacer tous les points de saturation d'un ordre de grandeur — et avec eux
toutes les cellules de régime. À un portique en rafale, les requêtes s'accumulent, et
un serveur réel les grouperait.

Les règles de MLPerf Inference~\cite{mlperf}, référence des bancs d'inférence,
répondent à ce problème en séparant explicitement des scénarios au lieu de laisser le
groupement libre : un scénario \emph{Server}, requête par requête, avec arrivées
aléatoires et contrainte de latence au 99\textsuperscript{e} percentile, et un
scénario \emph{Offline}, où toutes les requêtes sont disponibles d'emblée et seul le
débit compte.

\begin{encadre}{Recommandation D12}
Mesurer le scénario \textbf{Server sans groupement} comme scénario principal, et le
scénario \textbf{Offline} comme scénario secondaire déclaré, avec une taille de lot
gelée. La raison est une propriété du modèle : les deux bras ayant la même intensité
opérationnelle, le groupement déplace leurs débits absolus sans changer, au premier
ordre, leur rapport. La question du rapport se tranche donc sans groupement ; la
question de la faisabilité absolue — l'appareil tient-il la charge du portique ? — a
besoin du plafond Offline. Le groupement dynamique, qui attend un court délai pour
réunir des requêtes, est une politique de déploiement : son délai d'attente devrait
lui-même être gelé, et il relève d'une phase ultérieure.
\end{encadre}

\subsection{Le processus d'arrivée}

Le contrat prévoit des arrivées \emph{déterministes} : une requête toutes les
$1/\lambda$ secondes, exactement. MLPerf, dans son scénario Server, utilise des
arrivées \emph{poissoniennes}, c'est-à-dire aléatoires et sans mémoire. La différence
n'est pas un détail de méthode. C'est un résultat classique de la théorie des files :
avec des arrivées et un service tous deux déterministes, aucune requête n'attend tant
que la charge reste sous la capacité ; avec des arrivées poissoniennes et un service
déterministe — une file notée M/D/1 —, l'attente moyenne vaut
\[
  W = \frac{\rho}{2(1-\rho)} \ \text{durées de service},
\]
où $\rho$ est le taux d'occupation de l'appareil.

\begin{figure}[h]
\centering
\begin{tikzpicture}
\begin{axis}[
  width=11cm,height=6cm,
  xmin=0,xmax=0.96,ymin=0,ymax=10.5,
  xlabel={taux d'occupation $\rho$},
  ylabel={attente moyenne (durées de service)},
  axis lines=left,
  xtick={0,0.5,0.8,0.9},
  xticklabels={0,{0,5},{0,8},{0,9}},
  every axis label/.style={font=\small},
  tick label style={font=\small},
  legend style={font=\footnotesize,at={(0.03,0.97)},anchor=north west,draw=none,fill=white},
]
\addplot[very thick,rouge,domain=0:0.95,samples=120] {x/(2*(1-x))};
\addlegendentry{arrivées poissoniennes (M/D/1)}
\addplot[very thick,encre,domain=0:0.95,samples=2] {0};
\addlegendentry{arrivées déterministes (D/D/1)}
\addplot[only marks,mark=*,mark size=2pt,rouge] coordinates {(0.5,0.5) (0.8,2) (0.9,4.5)};
\node[font=\footnotesize,rouge,anchor=west] at (axis cs:0.5,1.1) {0,5};
\node[font=\footnotesize,rouge,anchor=east] at (axis cs:0.79,2.4) {2};
\node[font=\footnotesize,rouge,anchor=east] at (axis cs:0.89,4.9) {4,5};
\end{axis}
\end{tikzpicture}
\caption{L'écart entre les deux hypothèses d'arrivée explose près de la saturation —
précisément la zone où se décide la publication des rapports.}
\label{fig:md1}
\end{figure}

Des arrivées de stade sont au moins aussi irrégulières que des arrivées
poissoniennes, souvent davantage avant le coup d'envoi. Le banc actuel est donc
reproductible, mais \emph{optimiste sur la latence de queue}.

\begin{encadre}{Recommandation D13}
Conserver les arrivées déterministes comme processus canonique — les modifier
toucherait un construit déjà accepté —, et \textbf{ajouter une variante poissonienne à
graine gelée}, aux mêmes débits moyens. Une graine fixée rend la variante aussi
reproductible que la version déterministe. Publier les deux ; leur écart est en
soi une information sur la sensibilité de l'appareil à l'irrégularité du flux. Les
afflux plus violents que Poisson restent couverts, en partie, par l'essai en rafale.
\end{encadre}

\subsection{Le passage et la caméra}

Le troisième oubli a été traité en section~\ref{sec:operationnel} : la conversion
entre personnes par heure et requêtes par seconde.

\begin{encadre}{Recommandation D14}
Retenir comme conversion canonique \textbf{une requête par passage}, issue de la
sélection de la meilleure image, et déclarer une sensibilité à $k$ détections par
passage. Le nombre de caméras par appareil découle des classes de déploiement. Les
valeurs de $k$ et du nombre de caméras relèvent du scénario.
\end{encadre}

\subsection{La question produit prioritaire : quantifier plutôt que réduire}

\begin{table}[h]
\centering\small
\begin{tabular}{@{}lr@{}}
\toprule
représentation & octets par vecteur \\
\midrule
512D en FP32 & \num{2048} \\
512D en FP16 & \num{1024} \\
\textbf{512D en INT8} & \textbf{512} \\
\textbf{128D en FP32} & \textbf{512} \\
\bottomrule
\end{tabular}
\caption{Une représentation 512D quantifiée sur 8 bits occupe exactement la même
mémoire qu'une représentation 128D en virgule flottante 32 bits.}
\label{tab:int8}
\end{table}

C'est sans doute le constat le plus lourd de conséquences du chapitre. Tout l'argument
de mémoire et de palier matériel en faveur de 128D peut être obtenu \emph{sans réduire
la dimension}, par quantification seule. Même charge utile, même palier. Au premier
ordre, même latence aussi : l'intensité opérationnelle de 512D-INT8 vaut 2 opérations
par octet, toujours sous le point de crête, et le temps reste fixé par les octets lus.

Le benchmark a raison de séparer les deux interventions : la quantification y est une
phase distincte, et c'est la bonne construction pour isoler l'effet de la dimension.
Mais la \emph{décision produit} ne porte pas sur \og 512D contre 128D \fg. Elle porte
sur \og quelle représentation de 512 octets préserve le mieux la performance
biométrique \fg. Si 512D-INT8 se dégrade moins que 128D-FP32 à mémoire égale, la
réduction de dimension perd sa raison d'être du côté de l'ingénierie. L'étude
biométrique doit donc comparer les routes à \emph{octets égaux}, et le tableau de
décision doit porter une colonne \og octets par vecteur \fg, pas seulement une colonne
\og dimension \fg.

% =====================================================================
\section{Cadre réglementaire}\label{sec:reglementaire}
% =====================================================================

\begin{alerte}{Portée de cette section}
Ce qui suit établit quelles règles existent. Leur application au programme dépend de
la juridiction, de l'opérateur et de la finalité du traitement. C'est une
qualification juridique qui relève d'un conseil ; ce document n'en tient pas lieu.
\end{alerte}

\paragraph{Sécurité photobiologique.} L'{\NoAutoSpacing IEC 62471:2006} fixe les limites d'exposition,
la méthode de mesure et la classification des risques pour les sources optiques
incohérentes, diodes électroluminescentes comprises et lasers exclus, entre 200 et
\num{3000}~nm~\cite{iec62471}. C'est une norme horizontale : elle attribue un groupe
de risque, mais les exigences qui en découlent relèvent des normes propres à chaque
produit~\cite{bentham}. La nomenclature de l'appareil doit donc porter le groupe de
risque de l'illuminateur infrarouge.

\paragraph{Données biométriques, si le déploiement relève de l'Union européenne.}
L'article 9 du RGPD interdit par principe le traitement de données biométriques aux
fins d'identifier une personne de manière unique, sauf exceptions de son
paragraphe~2~\cite{gdpr}.

\paragraph{Intelligence artificielle en espace public, si le déploiement relève de
l'Union européenne.} Le règlement européen sur l'IA interdit en principe
l'identification biométrique à distance en temps réel dans les espaces accessibles au
public à des fins répressives (article 5, paragraphe 1, point h), sauf exceptions
étroites~\cite{aiact}. Les systèmes d'identification biométrique qui ne sont pas
interdits sont classés à haut risque.

\begin{alerte}{Conséquence pour le programme}
Les deux branches de l'analyse ne sont pas seulement deux familles de métriques : elles
peuvent relever de deux régimes juridiques. Un contrôle d'accès de porteurs de billets
et une liste de surveillance opérée à des fins répressives ne se qualifient pas de la
même manière. Dans l'Union européenne, la seconde peut être interdite par défaut,
quelle que soit la qualité de l'ingénierie. La juridiction et l'opérateur de chaque
branche sont des paramètres de scénario à déclarer \emph{avant} de dimensionner la
branche correspondante.
\end{alerte}

% =====================================================================
\section{Les quatorze décisions et l'ordre de gel}\label{sec:decisions}
% =====================================================================

Toutes les décisions ci-dessous se prennent sans regarder la moindre donnée. Chacune
porte un défaut argumenté, que le responsable du programme accepte ou remplace ; ce qui
compte est que la valeur soit gelée avant que quiconque voie la machine.

\begin{small}\renewcommand{\arraystretch}{1.28}
\begin{longtable}{@{}>{\bfseries}l>{\raggedright}p{3.6cm}>{\raggedright}p{6.6cm}l@{}}
\toprule
\normalfont id & objet & défaut proposé & statut \\
\midrule
\endhead
D1 & affectation de régime de saturation & rapport débit complété / offert : $\geq 0{,}98$ non saturé, $<0{,}90$ saturé ; départage par la pente de la file ; l'indéterminé compte comme régimes différents & défaut \\
D2 & débit dérivé obligatoire & charge utile × débit complété ÷ bande passante mesurée ; au-delà de 1, preuve de résidence en cache & défaut \\
D3 & cache effectif & somme des tranches de LLC des grappes où tourne le banc & défaut \\
D4 & caractérisation de bande passante & une fois, en début de session, suivie d'au moins 300~s de repos ; jamais entre deux blocs & défaut \\
D5 & comptes attendus et réalisés & publiés ensemble, avec suppression \og par conception \fg{} ou \og par déficit \fg & défaut \\
D6 & échelle de paliers & au moins deux paliers encadrant la faisabilité de 512D sur G2, chacun avec sa bande passante & à déclarer \\
D7 & marge de sécurité RAM & le plus grand de \SI{20}{\percent} de la RAM ou 512~Mo libres au pic & défaut \\
D8 & profils de service & deux ou trois profils horaires nommés & à déclarer \\
D9 & régime thermique & même construction que D1, avec repli si la puce n'expose pas de compteurs & défaut \\
D10 & banc thermiquement représentatif & mesure dans la configuration de déploiement la plus défavorable, sinon étiquette de meilleur cas & défaut \\
D11 & classes de déploiement & portique fixe et point mobile : alimentation, plafond, ambiant, boîtier, veille & à déclarer \\
D12 & groupement des requêtes & Server sans groupement en principal, Offline en secondaire & à trancher \\
D13 & processus d'arrivée & déterministe canonique, plus une variante poissonienne à graine gelée & à trancher \\
D14 & requêtes par passage & une par passage en canonique, sensibilité à $k$ détections & à déclarer \\
\bottomrule
\end{longtable}
\end{small}

\begin{figure}[h]
\centering
\begin{tikzpicture}[
  font=\sffamily\small,
  etape/.style={draw=encre,fill=bleupale,minimum width=9.4cm,minimum height=0.72cm,align=flush left,text width=9.1cm,inner xsep=6pt},
  cle/.style={etape,fill=ambrepale,draw=ambre},
  fl/.style={-{Stealth[length=2.5mm]},thick,encre},
  node distance=0.32cm]
\node[etape] (e1) {\textbf{1} \quad Geler la logique des gardes — D1 à D14};
\node[etape,below=of e1] (e2) {\textbf{2} \quad Revue indépendante de la spécification des gardes};
\node[etape,below=of e2] (e3) {\textbf{3} \quad Matérialiser la plateforme, lier les valeurs physiques};
\node[etape,below=of e3] (e4) {\textbf{4} \quad Geler l'enveloppe opérationnelle et les seuils experts};
\node[etape,below=of e4] (e5) {\textbf{5} \quad Exécuter la Phase A};
\node[cle,below=of e5] (e6) {\textbf{6} \quad Ouvrir SCREEN};
\node[etape,below=of e6] (e7) {\textbf{7} \quad Croiser biométrie, ingénierie et conséquences};
\foreach \a/\b in {e1/e2,e2/e3,e3/e4,e4/e5,e5/e6,e6/e7} \draw[fl] (\a) -- (\b);
\draw[decorate,decoration={brace,amplitude=5pt},gris,thick] ($(e3.north east)+(0.15,0)$) -- ($(e4.south east)+(0.15,0)$)
  node[midway,right=8pt,align=left,gris,font=\sffamily\footnotesize] {3 et 4 peuvent\\se faire en parallèle,\\4 jamais après 6};
\node[right=0.35cm of e6,align=left,ambre,font=\sffamily\footnotesize] {après la machine :\\règle n°~9};
\end{tikzpicture}
\caption{Ordre de gel. Le seul réordonnancement qu'impose la discussion produit est
l'échange de SCREEN et de la matérialisation.}
\label{fig:gel}
\end{figure}

Le premier lot de gardes à implémenter se choisit par la gravité de la conclusion
fausse, non par la facilité du test : l'absence de swap, l'isolement des bras, la
sémantique de recherche exacte, l'échauffement et l'alternance des bras, le
générateur hors appareil et en boucle ouverte, le plan de mesure énergétique, puis les
deux classifieurs de régime, de saturation et thermique.

Deux principes gouvernent leur écriture. Chaque garde sépare sa \emph{logique} —
formule, seuil, comportement en cas d'échec —, gelée maintenant, de sa
\emph{matérialisation} — compteur système, modèle de wattmètre —, qui viendra avec la
machine ; dans l'autre sens, on écrirait des assertions qui décrivent ce que la
plateforme fait au lieu de contraindre ce qu'elle doit faire. Et chaque garde conserve
la \emph{valeur observée}, pas seulement son verdict : sans cela, un relecteur ne peut
pas distinguer \og vérifié et correct \fg{} de \og n'a pas tourné \fg.

% =====================================================================
\section{Questions ouvertes}\label{sec:questions}
% =====================================================================

Aucune des questions suivantes ne se tranche en regardant des données. Toutes
relèvent du responsable du programme, et toutes doivent l'être avant le gel des gardes.

\paragraph{Valeurs de scénario.} L'échelle des paliers matériels candidats et leur
ordre d'essai (D6). Les profils de service heure par heure (D8). Les classes de
déploiement : pour chacune, l'alimentation, le plafond de puissance, la plage de
température ambiante, le type de boîtier, et si l'appareil peut dormir entre deux
passages (D11). La politique de sélection d'images et le nombre de caméras par appareil
(D14). Pour l'enveloppe opérationnelle : le nombre de passages attendus, la densité
d'appareils autour des sites, la taille de la liste de surveillance, et le point de
fonctionnement visé pour chaque branche.

\paragraph{Gouvernance.} Qui tient chacun des trois rôles de l'élicitation — détenteur
des résultats SCREEN, constructeur de l'enveloppe, expert — sachant qu'aucune personne
ne peut en tenir deux. La juridiction et l'opérateur de chaque branche, qui décident
du régime juridique applicable. La qualification juridique elle-même, par un conseil.

\paragraph{Défauts à accepter ou remplacer.} D1 à D5, D7, D9 et D10 portent un défaut
argumenté. D12 et D13 portent une recommandation qui demande un arbitrage explicite.

\paragraph{Sources à terminer.} Lire le détail du cache L3 dans la fiche NVIDIA
DS-11105-001. Consulter directement la fiche Intel du N100, non atteinte. Remplacer la
source secondaire sur les batteries par la fiche de la cellule effectivement retenue.
Vérifier la série de normes ISO/IEC 19795 sur l'évaluation des performances
biométriques avant de l'utiliser pour justifier une marge.

\paragraph{Question produit.} Faut-il introduire une route 512D quantifiée sur 8 bits
dans la comparaison biométrique future, pour comparer les représentations à octets
égaux ? Ce chapitre le recommande. Cette route relèverait d'une nouvelle identité de
protocole, sans toucher à Study 1B.

% =====================================================================
\section{Provenance des valeurs}\label{sec:provenance}
% =====================================================================

Chaque valeur factuelle de ce chapitre relève de l'une des natures suivantes. Une
conclusion ne peut reposer que sur les quatre premières.

\begin{description}[style=nextline,leftmargin=2.2cm,font=\sffamily\small\bfseries]
\item[primaire] fiche ou documentation du fabricant, texte normatif ou réglementaire
\item[secondaire] reprise par un tiers : presse technique, distributeur, analyse
\item[dérivée] calculée à partir d'une valeur documentée, formule donnée
\item[mesurée] mesurée sur l'appareil, méthode et date à l'appui — aucune à ce jour
\item[supposée] hypothèse de travail, interdite comme base d'une conclusion
\end{description}

Un résumé automatique d'une source primaire n'est pas une source, ni les réponses d'un
assistant, y compris celles qui ont servi à rédiger ce chapitre.

\begin{small}
\begin{longtable}{@{}>{\raggedright}p{5.7cm}>{\raggedright}p{4.1cm}ll@{}}
\toprule
fait & valeur & nature & source \\
\midrule
\endhead
Raspberry Pi 5 — cache L3 partagé & 2~Mo & primaire & \cite{rpidoc} \\
Raspberry Pi 5 — bande passante mémoire & jusqu'à 17~Go/s & primaire & \cite{rpidoc} \\
Raspberry Pi 5 — seuils de throttling & \SI{80}{\celsius} et \SI{85}{\celsius} & primaire & \cite{rpiheat} \\
Raspberry Pi 5 — sans refroidissement, charge soutenue & throttling permanent & primaire & \cite{rpiheat} \\
Raspberry Pi 5 — puissance sous charge & environ 12~W & secondaire & \cite{nbrpi} \\
Jetson Orin Nano — bande passante de crête & 34~Go/s (4~Go), 68~Go/s (8~Go) & primaire & \cite{nvds} \\
Jetson Orin Nano — canaux vidéo virtuels & jusqu'à 16 & primaire & \cite{nvds} \\
Jetson Orin Nano — exigence de conception thermique & SoC sous sa température max & primaire & \cite{nvtdg} \\
Jetson Orin Nano — cache L3 & 2~Mo par grappe & secondaire & \cite{ridgerun} \\
Intel N100 — cache L3, mémoire & 6~Mo, canal unique DDR5-4800 & secondaire & \cite{nbcheck} \\
Intel N100 — bande passante de crête & \num{38,4}~Go/s, soit \num{4800}~MT/s × 8 octets & dérivée & calcul \\
PoE 802.3af et 802.3at — puissance à l'appareil & \num{12,95}~W et \num{25,5}~W & secondaire & \cite{poe} \\
Cellule lithium-ion — plages de charge et de décharge & 0 à \SI{45}{\celsius} ; $-20$ à \SI{60}{\celsius} & secondaire & \cite{bu410} \\
Caméra extérieure chauffée — pointe hivernale & jusqu'à 32~W & secondaire & \cite{ampcom} \\
IEC 62471 — portée & 200 à \num{3000}~nm, lasers exclus & catalogue normatif & \cite{iec62471} \\
RGPD, article 9 & interdiction de principe & texte réglementaire & \cite{gdpr} \\
AI Act, article 5(1)(h) & interdiction de principe & analyse juridique & \cite{aiact} \\
Modèle Roofline & $\min(\text{crête},\ \text{BP} \times \text{IO})$ & primaire & \cite{roofline} \\
Scénarios MLPerf Inference & Server, Offline et autres & primaire & \cite{mlperf} \\
Méthode STREAM & noyau Triad de référence & secondaire & \cite{stream} \\
Efficacité de lecture continue & \SI{70}{\percent} & \textbf{supposée} & — \\
Surcoût système et exécution & 1~Go & \textbf{supposée} & — \\
\bottomrule
\end{longtable}
\end{small}

\appendix

% =====================================================================
\section{Glossaire}\label{sec:glossaire}
% =====================================================================

\begin{description}[style=unboxed,leftmargin=0pt,font=\sffamily\small\bfseries,itemsep=3pt]
\item[AI Act] Règlement européen sur l'intelligence artificielle.
\item[BLAS] \emph{Basic Linear Algebra Subprograms} — bibliothèque standard de calcul
  matriciel, qui choisit elle-même ses micro-noyaux selon la forme du calcul.
\item[Bras] Configuration de mesure complète sous une seule condition ; la Phase A en
  compte deux, 512D et 128D.
\item[D/D/1, M/D/1] Notations de files d'attente à un serveur : arrivées déterministes
  ou poissoniennes (D ou M), service déterministe (D).
\item[DRAM] Mémoire vive dynamique, la mémoire principale de l'appareil.
\item[FMR, FNMR] En vérification \unun : taux de fausses acceptations d'imposteurs, et
  de faux rejets de personnes légitimes.
\item[FP32, FP16, INT8] Formats numériques : virgule flottante sur 32 ou 16 bits,
  entier sur 8 bits. 4, 2 ou 1 octet par composante.
\item[FPIR, FNIR] En identification \unN : taux de fausses alertes pour des personnes
  absentes de la liste, et de personnes de la liste passées sans alerte.
\item[Garde] Mécanisme qui vérifie qu'une clause est tenue et bloque sinon : test de
  contrat, assertion d'exécution ou contrôle humain.
\item[Intensité opérationnelle] Nombre d'opérations effectuées par octet lu en mémoire.
\item[IR, NIR] Infrarouge, proche infrarouge.
\item[LFW] \emph{Labeled Faces in the Wild}, jeu académique de visages, \num{5749}
  identités.
\item[LLC] \emph{Last level cache} — dernier niveau de cache du processeur, le plus
  grand et le plus lent.
\item[LPDDR] Famille de mémoires vives basse consommation pour appareils embarqués.
\item[MLPerf] Suite de bancs d'essai d'inférence de MLCommons, dont les règles
  définissent des scénarios de charge.
\item[NI, $\delta_{\mathrm{NI}}$] Non-infériorité ; marge de non-infériorité.
\item[PoE] \emph{Power over Ethernet} — alimentation d'un appareil par son câble réseau.
\item[Point de fonctionnement] Réglage du seuil de décision qui fixe un taux d'erreur
  pour mesurer l'autre.
\item[RGPD] Règlement général sur la protection des données (GDPR en anglais).
\item[Roofline] Modèle de performance : le minimum entre crête de calcul et bande
  passante × intensité opérationnelle.
\item[SCREEN, TEST] Étapes de Study 1B : exploration biométrique non décisionnelle,
  puis qualification. Toutes deux scellées.
\item[SoC] \emph{System on chip} — puce réunissant processeur, mémoire cache et
  contrôleurs.
\item[STREAM] Banc de référence pour la bande passante mémoire soutenue.
\item[Swap] Débordement de la mémoire vive sur le stockage, très lent.
\item[Throttling] Réduction automatique de fréquence du processeur au-delà d'un seuil
  de température.
\item[Working set] Données effectivement touchées par un programme pendant son travail.
\end{description}

% =====================================================================
\begin{thebibliography}{99}
% =====================================================================
\small
\bibitem{roofline} S. Williams, A. Waterman, D. Patterson. \emph{Roofline: An
  Insightful Visual Performance Model for Multicore Architectures}. Communications of
  the ACM 52(4), avril 2009, p.~65--76. DOI 10.1145/1498765.1498785. Version libre :
  rapport technique UC Berkeley UCB/EECS-2008-134.
\bibitem{hp} J. L. Hennessy, D. A. Patterson. \emph{Computer Architecture: A
  Quantitative Approach}, 6\textsuperscript{e} édition. Morgan Kaufmann, 2017.
\bibitem{stream} STREAM, banc de bande passante mémoire soutenue. Documentation de
  référence, dont Microsoft Learn, \emph{STREAM benchmark}.
\bibitem{mlperf} MLCommons. \emph{MLPerf Inference Rules}, dépôt
  \texttt{mlperf/inference\_policies}, fichier \texttt{inference\_rules.adoc}.
\bibitem{rpidoc} Raspberry Pi Ltd. Documentation matérielle, processeur BCM2712.
  \url{https://www.raspberrypi.com/documentation/hardware/raspberrypi/bcm2711/}
\bibitem{rpiheat} Raspberry Pi Ltd. \emph{Heating and cooling Raspberry Pi 5}.
  \url{https://www.raspberrypi.com/news/heating-and-cooling-raspberry-pi-5/}
\bibitem{nbrpi} Notebookcheck. \emph{Broadcom BCM2712}.
\bibitem{nvds} NVIDIA. \emph{Jetson Orin Nano Series Data Sheet}, DS-11105-001.
\bibitem{nvtdg} NVIDIA. \emph{Jetson Orin NX Series and Jetson Orin Nano Series
  Modules Thermal Design Guide}, TDG-11127-001.
\bibitem{ridgerun} RidgeRun. \emph{NVIDIA Jetson Orin Nano — SoM Overview}.
\bibitem{nbcheck} Notebookcheck, TechPowerUp. Fiches du processeur Intel N100.
\bibitem{poe} Spécifications IEEE 802.3af et 802.3at, telles que rapportées par
  plusieurs sources concordantes (Netgear, Schneider Electric, Faceofit). Norme IEEE
  payante, non consultée.
\bibitem{bu410} Battery University. \emph{BU-410: Charging at High and Low
  Temperatures}.
  \url{https://batteryuniversity.com/article/bu-410-charging-at-high-and-low-temperatures}
\bibitem{sunlith} SunLith Energy. \emph{Charging temperature for batteries}.
\bibitem{ampcom} AMPCOM. \emph{PoE Standards: 802.3af, 802.3at \& 802.3bt Complete
  Glossary}.
\bibitem{iec62471} {\NoAutoSpacing IEC 62471:2006}, \emph{Photobiological safety of lamps and lamp
  systems}. Notice de catalogue, Conseil canadien des normes et SIS.
\bibitem{bentham} Bentham Instruments. \emph{Photobiological Safety of Lamps and Lamp
  Systems}.
\bibitem{gdpr} Règlement (UE) 2016/679, article 9. Texte reproduit par gdpr-text.com.
\bibitem{aiact} Règlement européen sur l'intelligence artificielle, article 5. Analyse :
  Verfassungsblog, \emph{AI Act and the Prohibition of Real-Time Biometric
  Identification}, décembre 2024.
\end{thebibliography}

\end{document}
